% !TeX root = ../main.tex
% chapters/background.tex

\chapter{Background}
\label{ch:background}

% ── Domain theory ─────────────────────────────────────────────────
\section{Domain Theory}
\label{sec:bg-domain-theory}

% Source files: structures/Order.v, structures/CPO.v
% Context: OrderTheory, CPOTheory _theory.md files

\subsection{Partial Orders and $\omega$-Chains}
\label{sec:bg-partial-orders}

% TODO: Preorders, partial orders, ω-chains, monotone functions
% Key formalization: Order.v (HasLe, IsPreorder, IsPartialOrder, chain)

\subsection{$\omega$-Complete Partial Orders}
\label{sec:bg-cpo}

% TODO: CPOs, suprema, Scott continuity, pointed CPOs
% Key formalization: CPO.v (HasSup, IsCPO, HasBottom, IsPointed)

\subsection{Fixed-Point Theory}
\label{sec:bg-fixedpoints}

% TODO: Kleene fixed-point theorem, Scott induction, fixp_mono
% Key formalization: theory/FixedPoints.v

\subsection{Scott Topology}
\label{sec:bg-scott-topology}

% TODO: Way-below relation, Scott-open sets, continuity equivalence
% Key formalization: theory/ScottTopology.v

% ── Enriched category theory ─────────────────────────────────────
\section{Enriched Category Theory}
\label{sec:bg-enriched}

% Source files: structures/Enriched.v, references: Kelly1982, MacLane1998

\subsection{Enriched Categories}
\label{sec:bg-enriched-cats}

% TODO: V-enriched categories, hom-objects, composition, identity
% Specialize to V = CPO

\subsection{Enriched Functors and Natural Transformations}
\label{sec:bg-enriched-functors}

% TODO: Locally continuous functors, enriched nat trans

% ── Denotational semantics ────────────────────────────────────────
\section{Denotational Semantics}
\label{sec:bg-denotational}

% References: BKV2009, Plotkin1977

\subsection{PCF: Programming Computable Functions}
\label{sec:bg-pcf}

% TODO: PCF syntax and operational semantics (informal)

\subsection{Soundness and Adequacy}
\label{sec:bg-adequacy}

% TODO: What soundness and adequacy mean, why they matter

% ── Proof assistants ──────────────────────────────────────────────
\section{Proof Assistants and \Rocq{}}
\label{sec:bg-rocq}

% References: HB2023

\subsection{The \Rocq{} Proof Assistant}
\label{sec:bg-rocq-assistant}

% TODO: Brief intro to Rocq (CIC, tactics, extraction)

\subsection{Hierarchy Builder}
\label{sec:bg-hb}

% TODO: HB mixin/structure/instance pattern, why it matters
